\section{Closing the uniform symplectic Frobenius catalogue}
\label{sec:CX}

\status{CX1--CX9 supply the literal \texttt{NumberSp4Alt}.  The
proof uses the precisely stated classical results below and the
already reconstructed polynomial coefficient modules and SP9.  It
has no retained bespoke catalogue, automorphy-counting, or
level-comparison premise.  An individual representation must still
satisfy the original \texttt{AdmRAlt} hypotheses; constructing its
\texttt{IsAutQA} witness is a different obligation.}

The useful target is a catalogue of four-dimensional
representations with weights
\[
 m=(e+1,e,1,0),\qquad w=e+1,\qquad e\ge2.                 \tag{CX.1}
\]
We prove exactly the catalogue used by SP9, with its labels chosen
before the varying coefficient prime, algebraic closure, complex
embedding, and representation family.  The main simplification is
to allow the auxiliary automorphic representation to have level
$M^7$.  Its original Galois representation continues to have level
$M$.  Since a fixed power of $M$ only changes a polynomial counting
exponent, the final source predicate permits this simplification.
We do not assert the source's stronger, unchanged-level
\texttt{LevelPNAlt} or its all-rank \texttt{CountQAlt}.

\subsection{Conventions and the classical theorems being used}

For a number field $F$, $G_F$ denotes its absolute Galois group in
the specified algebraic closure.  At a finite place $v$, write
$q_v$ for the residue-field cardinality.  Reciprocity sends a
uniformizer to \emph{geometric} Frobenius.  Thus the cyclotomic
character has value $q_v^{-1}$ there.  The source's
\texttt{FrobEig} uses this geometric convention, whereas
\texttt{NumberSp4Alt} evaluates an \emph{arithmetic} Frobenius.
If $B_v(\pi)$ is the usual unramified Satake multiset for
$\mathrm{GL}_4$, put
\[
 A_v(\pi)=q_v^{3/2}B_v(\pi),\qquad
 \beta_v(\pi)=\{a^{-1}:a\in A_v(\pi)\}.                \tag{CX.2}
\]
The first multiset is the geometric Galois multiset; the second is
the arithmetic multiset.  The factor $q_v^{3/2}$ is exactly the
source's \texttt{algSat}, equivalently local Langlands for
$\pi_v\otimes|\det|^{-3/2}$.

Here are the classical inputs, in the forms actually used.  Their
primary references and exact source-interface matching are also
recorded in the component's classical-input registry.

\begin{enumerate}
\item A regular algebraic cuspidal representation of
$\mathrm{GL}_n$ has nonzero relative Lie algebra cohomology with
the irreducible algebraic coefficient of the corresponding
infinitesimal character.  Its finite part has a number-field
model, and every automorphism of $\mathbb C$ produces a cuspidal
conjugate with the permuted infinity type.  These are Clozel's
cohomological and rationality theorems, specifically Lemma~3.14
and Theorem~3.13 of \emph{Motifs et formes automorphes}.
\item Cuspidal cohomology injects into ordinary cohomology with
the same coefficient system.  We use the arithmetic case of
Borel--Labesse--Schwermer, Theorem~5.2.  The Borel--Serre
compactification gives a finite CW model for each fixed
torsion-free arithmetic quotient.  These assertions alone do not
provide a uniform counting theorem; CX1 proves the required
uniform bound from the degree of a congruence covering.
\item For a regular algebraic cuspidal polarized pair
$(\pi,\chi)$ over a CM field, and every prime $a$ and isomorphism
$\overline{\mathbb Q}_a\simeq\mathbb C$, the associated
semisimple $a$-adic representation exists.  At places away from
$a$, it satisfies full Frobenius-semisimplified local-global
compatibility and purity.  Its Hodge--Tate numbers are
$\lambda_{\tau,i}+n-i$ when the coefficient has highest weight
$\lambda_\tau$; a spherical place over $a$ gives a crystalline
representation.  This is Barnet-Lamb--Gee--Geraghty--Taylor,
\emph{Potential automorphy and change of weight},
Theorem~2.1.1.  At an arbitrary place over $a$ we use only the
\emph{semisimplified} compatibility of their
\emph{Local-global compatibility for $l=p$, II}, Theorem~A.
No Shin-regularity or monodromy equality at $a$ is assumed.
\item Every de Rham representation is potentially semistable
(Berger, \emph{Repr\'esentations $p$-adiques et \'equations
diff\'erentielles}, Theorem~0.6, equivalently Theorem~5.19).
The potentially semistable comparison functor is exact
on potentially semistable representations; a crystalline
representation gives a weakly admissible filtered Frobenius
module, with the same Hodge--Tate numbers.  We use the classical
Fontaine comparison theorem in the form recalled by
Colmez--Fontaine, Proposition~4.2 and the corollary to
Theorem~4.3.  CX3 spells out the slope inequality needed here.
\item Local Langlands for general linear groups preserves
conductors, is compatible with characters and unramified
parameters, and every local component of a cuspidal
$\mathrm{GL}_n$ representation is generic.  The generic newvector
theorem of Jacquet--Piatetski-Shapiro--Shalika, Theorem~5.1,
then gives a vector fixed by the usual $K_1(\mathfrak p^a)$ at
conductor exponent $a$.  We use only the resulting invariance
under the smaller principal congruence subgroup.  Conductor
compatibility is the $\varepsilon$-factor part of local
Langlands; Scholze's Theorem~14.1 is a primary reference.
\item We use global and local class field theory, Chebotarev,
and strong multiplicity one for cuspidal $\mathrm{GL}_n$.
CX4 writes the cyclotomic character, polarization twist and
comparison of the original representations explicitly.  None of
these inputs asserts automorphy of an arbitrary Galois
representation.
\end{enumerate}

These theorems have their usual characteristic-zero hypotheses.
In particular, purity below applies only after a polarized
regular algebraic cuspidal representation has actually been
constructed.  The sole weak automorphy witness used for that
construction is already a field of \texttt{AdmRAlt}.

\subsection{A uniform CM catalogue, including its arithmetic properties}

\begin{proposition}[CX1: uniform counting by fixed cells and polynomial coefficients]
\label{claim:CX1}
For each fixed imaginary quadratic field $K$ there is an integer
$c_K\ge1$ such that, for every $e\ge2$ and $L\ge1$, the number
of isomorphism classes of cuspidal representations $\pi$ of
$\mathrm{GL}_4(\mathbb A_K)$ with
\[
 \operatorname{Inf}_\tau(\pi)=\{3/2+m_i:0\le i<4\}
 \quad(\tau:K\hookrightarrow\mathbb C),\qquad
 \pi_f^{K_K(L)}\ne0                                         \tag{CX.3}
\]
is at most $c_K L^{32}(e+2)^{32}$.  Here $K_K(L)$ is the
literal principal congruence subgroup modulo the integer $L$.
The same bound holds after restricting to polarized pairs.
\end{proposition}

\begin{proof}
The required coefficient has, at both embeddings, highest weight
\[
 \lambda=(-3,-3,-e-1,-e-1).
                                                               \tag{CX.4}
\]
Indeed its contragredient has infinitesimal character
$\{3/2+m_i\}$.  More explicitly, if indices are $1,\ldots,4$,
then $\lambda_i+4-i=-m_{5-i}$; reversing these entries gives
the prescribed character of the dual coefficient.

We can use the polynomial realization already proved in
HP1--HP4.  Set $H=e+1$.  The source's polynomial highest weight
is
\[
 \nu_i=i+H-m_{5-i}=(e+2,e+2,4,4)_i,
 \qquad \lambda=\nu-(e+5)(1,1,1,1).
                                                               \tag{CX.5}
\]
The polynomial module is contained in the polynomials of total
degree $2e+12$ in sixteen matrix entries.  Counting all possible
exponent tuples, even those of smaller degree, bounds its
dimension by $(2e+13)^{16}$.  Tensoring the two embedding factors
and applying the indicated determinant twists therefore gives a
coefficient system $E_e$ of rank
\[
             d_e\le(2e+13)^{32}\le5^{32}(e+2)^{32}.
                                                               \tag{CX.6}
\]
This argument uses the actual highest weight; it does not count
all representations of a merely comparable dimension.

Consider the adelic locally symmetric space with level $K_K(3)$,
with the positive real central factor removed in the usual
cohomological convention.  It has finitely many arithmetic
components.  For $\mathrm{GL}_4$ their finiteness also follows
directly by classifying the underlying rank-four
$\mathcal O_K$-lattices: each is isomorphic to
$\mathcal O_K^3\oplus\mathfrak a$, and the ideal class group is
finite.  Passing from level one to level three only introduces a
finite index.  The level-three groups are torsion free.  To see
the last assertion, use an integral basis of each lattice to
embed its stabilizer in $\mathrm{GL}_8(\mathbb Z)$; an element
congruent to $1$ modulo $3$ cannot have finite order unless it is
$1$.  If a nonidentity element of prime order were
$1+3^sA$, with $s\ge1$ maximal, the binomial expansion of its
prime power, divided by the first nonzero power of $3$, gives a
nonzero reduction modulo $3$, a contradiction.  The same
argument for a prime-order power excludes every finite order.

Choose finite CW models from Borel--Serre for all these fixed
components, and let $b_K$ be their total number of cells.  The
level $K_K(3L)$ space is a finite cover.  Its degree is bounded by
\[
 [K_K(1):K_K(3L)]
 \le\bigl|M_4(\mathcal O_K/3L\mathcal O_K)\bigr|
 =(3L)^{32}.                                                  \tag{CX.7}
\]
Indeed reduction modulo $3L$ identifies this quotient with a
subgroup of the invertible matrices in the displayed finite
matrix ring.  Each base cell has at most this many lifts.  A
rank-$d_e$ local system contributes $d_e$ dimensions per cell to
the cellular cochain complex, independently of its attaching
maps or monodromy.  Hence the sum of all cohomology dimensions
is at most $b_K(3L)^{32}d_e$.

Every $\pi$ in (CX.3) also has $K_K(3L)$-fixed vectors.
Clozel's cohomological theorem and cuspidal-cohomology injection
put a nonzero summand for $\pi_f^{K_K(3L)}$ in the cohomology
just bounded.  Distinct cuspidal isomorphism classes give
distinct summands; equivalently their full unramified Hecke
eigensystems are distinct by strong multiplicity one.
Consequently any finite set of distinct classes has at most
$b_K3^{32}5^{32}L^{32}(e+2)^{32}$ members.  The whole set is
therefore finite with that bound: otherwise one could select a
finite subset of larger cardinality.  Enlarging the fixed
integer constant proves the proposition.  In particular the
constant is chosen before $L$, $e$, or any Frobenius prime.
\end{proof}

\begin{proposition}[CX2: the full labels and their conjugation action]
\label{claim:CX2}
Let $\mathcal S_K(e,L)$ be the finite set of isomorphism classes
in (CX.3) that admit a polarized pair.  Every
$\gamma\in\operatorname{Aut}(\mathbb C)$ permutes this set, and
at every $v\nmid L$,
\[
 A_v({}^{\gamma}\pi)=\gamma A_v(\pi),\qquad
 \beta_v({}^{\gamma}\pi)=\gamma\beta_v(\pi).
                                                               \tag{CX.8}
\]
All these multisets consist of algebraic numbers.  The labels
remain the full global isomorphism classes; classes with equal
eigenvalues at one prime are not identified.
\end{proposition}

\begin{proof}
Clozel's conjugation theorem gives ${}^{\gamma}\pi$, whose
finite part is the scalar conjugate of $\pi_f$.  Fixed vectors
for the actual compact subgroup are carried to fixed vectors by
scalar conjugation.  The infinity type is unchanged because
(CX.3) is the same at both embeddings.  To check polarization,
write
\[
                 \pi^c\simeq\pi^\vee\otimes(\chi\circ N_{K/\mathbb Q}),
 \qquad \chi_\infty(-1)=1.                                   \tag{CX.9}
\]
Here $c$ is complex conjugation of $K$, and $\chi$ is an
algebraic Hecke character of $\mathbb Q$.  Such a character is a
finite Dirichlet character times an integral norm power.
Conjugate the finite character and keep the integral power.
This produces $\chi_\gamma$ with the same sign.  Scalar
conjugation commutes with the contragredient, the permutation of
places induced by $c$, and this character twist.  Thus (CX.9)
holds at almost every place for ${}^{\gamma}\pi$ and
$\chi_\gamma$.  Strong multiplicity one makes it a global
identity.  This proves that the polarized subset is stable.
Applying $\gamma^{-1}$ shows that the action is a permutation.

There is a useful normalization check at an unramified place.
Normalize Haar measure so that $\mathrm{GL}_4(\mathcal O_{K_v})$
has volume one, and let $T_{v,j}$ be the spherical Hecke
operator for the diagonal element with $j$ uniformizers.  Its
eigenvalue satisfies
\[
 e_j(A_v(\pi))=q_v^{j(j-1)/2}T_{v,j}(\pi),\qquad 0\le j\le4.
                                                               \tag{CX.10}
\]
The exponent on the right is an integer.  The Hecke operator is
defined over $\mathbb Q$, so its eigenvalue is conjugated by
$\gamma$.  Equality of the elementary symmetric polynomials
gives the first equality in (CX.8), including the potentially
misleading square-root factor in (CX.2).  Inversion gives the
second equality.  Clozel's number-field model for $\pi_f$
makes the eigenvalues of $T_{v,j}$ algebraic, so (CX.10) makes
the entries of $A_v$ and $\beta_v$ algebraic as well.
\end{proof}

\begin{proposition}[CX3: integral arithmetic eigenvalues and their size]
\label{claim:CX3}
Let $\pi\in\mathcal S_K(e,L)$, let $p\nmid L\operatorname{disc}K$
be prime, and let $v\mid p$, with $q=q_v$.  Every
$b\in\beta_v(\pi)$ is an algebraic integer, as is $q^{e+1}b^{-1}$.
Moreover
\[
                 |b|=q^{(e+1)/2}.                            \tag{CX.11}
\]
The assertions hold for every complex conjugate with the same
$K,e,L,p,v$.  In particular this supplies the source's full
\texttt{int} and \texttt{intInv} conditions for the weights
(CX.1), and supplies its weaker \texttt{Win} condition whenever
$q>1$.
\end{proposition}

\begin{proof}
Apply the polarized automorphic Galois theorem with any
auxiliary prime $a$.  From (CX.4), at every embedding its
Hodge--Tate numbers are
\[
                     0,-1,-e,-e-1.                          \tag{CX.12}
\]
The weight of its pure Weil--Deligne representation at
$v\nmid a$ is therefore $-e-1$: the theorem pairs the Hodge
numbers at the two conjugate embeddings by $h\mapsto w_0-h$,
and (CX.12) forces $w_0=-e-1$.  Its geometric Frobenius
multiset is $A_v(\pi)$ by local-global compatibility.  Purity
gives $|a_v|=q^{-(e+1)/2}$ for those eigenvalues, which proves
(CX.11) after inversion.  CX2 gives the same argument for every
complex conjugate.

We check integrality at \emph{every} finite place of the number
field generated by $b$.  First choose a prime $a\ne p$ and an
embedding of that number field into
$\overline{\mathbb Q}_a$.  Extend its inverse to the complex
coefficient isomorphism used in the automorphic Galois theorem.
Its representation can be written over a finite extension of
$\mathbb Q_a$.  For completeness, the finite-coefficient
capture needed here follows from the same Baire argument as
RV4: enumerate the finite extensions of $\mathbb Q_a$ inside
its algebraic closure; the closed subgroups on which all matrix
entries lie in one such field cover the compact profinite
domain.  One has interior and is therefore open.  Adjoining
the finitely many matrices for its cosets gives one finite
coefficient field for the whole representation.

The compact image preserves a lattice.  Indeed the sum of all
translates of one lattice is bounded above by a fixed
fractional lattice, contains the original lattice, and is a
finitely generated module over the coefficient DVR.  It is
therefore a full stable lattice.  Frobenius and its inverse
act invertibly on it.  Their eigenvalues, including the chosen
image of $b$, are $a$-adic units.

For the prime $p$ itself, use the associated $p$-adic
representation.  The local factor $\pi_v$ is spherical, so
the representation is crystalline at $v$.  Let $f$ be the
residue degree, so $q=p^f$.  The eigenvalues of $\varphi^f$ in
its crystalline module are the images of $A_v(\pi)$.  Weak
admissibility and (CX.12) put every normalized Newton slope in
$[-e-1,0]$.  Here is the needed elementary implication.  A
lowest-slope Frobenius subobject, with its induced filtration,
has $t_N\ge t_H$, and every filtration jump is at least
$-e-1$; its slope cannot be below $-e-1$.  A highest-slope
quotient has $t_N\le t_H$, by subtracting the subobject
inequality from the equality for the whole module.  Every
filtration jump of the quotient is at most zero, so its slope
cannot exceed zero.  The slope subobject or quotient can be
formed after a finite coefficient extension by grouping the
generalized eigenvectors of $\varphi^f$ with the indicated
valuation, together with their cyclic Frobenius translates.
This proves the bounds without assuming a splitting of the
original crystalline representation.  The fact that $p$ is
unramified in $K$ makes $\varphi^f$ precisely the Frobenius
with residue cardinality $q$; the same argument is applied to
each coefficient embedding.

Consequently, for a geometric eigenvalue $a_v=b^{-1}$,
\[
             -f(e+1)\le v_p(a_v)\le0.
                                                               \tag{CX.13}
\]
Thus $b$ and $q^{e+1}b^{-1}$ have nonnegative valuation at
every place over $p$, and they have nonnegative valuation at
all other finite places by the preceding unit argument.
Since they are algebraic, they are algebraic integers.
For the source's \texttt{int} condition, $-a\le m_i$ for all
$i$ forces $a\ge0$, because $\min m_i=0$.  For
\texttt{intInv}, $m_i\le a$ forces $a\ge e+1$.
Multiplication by a further nonnegative integral power of $q$
therefore gives exactly both universally quantified conditions.
Finally (CX.11) gives the correct product modulus and places
every entry strictly between $q^{(e+1)/2}/\sqrt q$ and
$\sqrt q\,q^{(e+1)/2}$.  This is \texttt{Win}.
\end{proof}
\subsection{Attaching the actual representation at a polynomial level}

\begin{proposition}[CX4: polarization and comparison in the original closure]
\label{claim:CX4}
Fix an odd prime $\ell$, an original algebraic closure
$\Omega/\mathbb Q$, an isomorphism
$\iota:\overline{\mathbb Q}_\ell\simeq\mathbb C$, and a
representation $r$ satisfying the literal
\texttt{AdmRAlt} at $(M,m,w)$ from (CX.1).  Suppose that an
imaginary quadratic subfield $K\subset\Omega$ has irreducible
restriction $r_K$.  The \texttt{IsAutQA} field supplies a
regular algebraic cuspidal representation $\pi/K$ weakly
attached to $r_K$.  This same $\pi$ admits a polarized pair,
and its associated $\ell$-adic representation is isomorphic to
the actual $r_K$ in the same coefficient field and closure.
Its infinity type is (CX.3), and at every place away from
$\ell$ its full Weil--Deligne parameter is the one of $r_K$.
\end{proposition}

\begin{proof}
Continuity is a field of the original \texttt{Level}.  Thus
the literal \texttt{IsAutQA} hypothesis applies to this very
quadratic restriction, with the given $\Omega$ and $\iota$.
It provides $\pi$ and equality of geometric Frobenius
polynomials at all but finitely many places.  It also provides
the additional \texttt{WitA} assertion, which we need not
strengthen or use for this argument.

Write the original multiplier as
\[
 \theta=\varepsilon_\ell^{e+1}\eta,
 \qquad \eta\text{ factors through }
       \operatorname{Gal}(\mathbb Q(\mu_N)/\mathbb Q)
 \quad\text{for some }N\ge1.                                 \tag{CX.14}
\]
This is exactly \texttt{IsCycWt}; factoring follows because its
\texttt{KillsMu} condition kills the fixing subgroup of all
$N$-th roots of unity.  The finite character is continuous.
Class field theory makes
$\theta\varepsilon_\ell^3$ the $\ell$-adic character of an
algebraic Hecke character $\chi$ of $\mathbb Q$.  In particular
at a good rational place, with geometric Frobenius $F_v$,
\[
             \chi(v)=\iota\theta(F_v)\,q_v^{-3}.
                                                               \tag{CX.15}
\]
One can construct it directly from the finite Dirichlet
character in (CX.14) and the integral norm power $e+4$.
This is the source's shift by $n-1=3$; using only $\theta$
would give the wrong automorphic twist.

Choose a lift in $G_{\mathbb Q}$ of complex conjugation on
$K$.  Conjugating $r_K$ by its matrix under $r$ identifies
$r_K^c$ with $r_K$.  The original nondegenerate bilinear
form in \texttt{DualTw} identifies $r_K$ with
$r_K^\vee\otimes\theta|_{G_K}$.  At a common unramified
place, put $a=q_v^{3/2}b$, where $a$ is a geometric Galois
eigenvalue and $b$ a Satake eigenvalue.  The duality relation
sends $a$ to $\theta(F_v)/a$.  After dividing by $q_v^{3/2}$
it sends $b$ to $\theta(F_v)q_v^{-3}/b$.  By (CX.15) these
are precisely the Satake parameters of
$\pi^\vee\otimes(\chi\circ N_{K/\mathbb Q})$.
Conjugation of places gives those of $\pi^c$.
Strong multiplicity one proves (CX.9).

If $\chi_\infty(-1)$ is not $+1$, multiply $\chi$ by the
quadratic character $\eta_K$ of $K/\mathbb Q$.
This changes the real sign and leaves
$\chi\circ N_{K/\mathbb Q}$ unchanged: local reciprocity
sends a norm into the subgroup fixing $K$, or equivalently the
quadratic character is trivial on these norms.  It also leaves
the restriction of the associated $\ell$-adic character to
$G_K$ unchanged.  Thus
$\varepsilon_\ell^{-3}\chi_\ell|_{G_K}=\theta|_{G_K}$
after this correction; no equality of the corrected characters
on all of $G_{\mathbb Q}$ is asserted.  We now have the literal polarized-pair sign
for rank four.  No polarization theorem about an arbitrary
nonautomorphic representation has been invoked.

Apply the polarized automorphic Galois theorem to this pair.
Its representation $R_{\ell,\iota}(\pi)$ is semisimple and
continuous.  The original $r_K$ is irreducible and continuous.
Their geometric Frobenius polynomials agree at almost every
place, by weak attachment and the theorem's normalization
(CX.2).  Both images lie in a finite coefficient extension by
the Baire argument recorded in CX3.  Chebotarev and continuity
make their traces agree on every element of $G_K$.
Equal characters of semisimple representations in
characteristic zero give an isomorphism.  For an elementary
form of this last step, take the finite-dimensional algebra
spanned by the image on their direct sum.  Its radical acts
trivially on this semisimple faithful module, so the algebra
is a product of matrix algebras.  Equality of traces on that
algebra, and then on its central idempotents, gives equal
multiplicities of every simple factor.  This proves
\[
                      r_K\simeq R_{\ell,\iota}(\pi).
                                                               \tag{CX.16}
\]

The infinity type is also forced, not chosen.  The original
flag at $\ell$ has one-dimensional quotients $\kappa_i$ with
inertial characters $\varepsilon_\ell^{m_i}\eta_i$, where
$\eta_i$ is finite.  Each quotient differs from that character
by an unramified character and has Hodge--Tate number $-m_i$.
Restriction to $K$ preserves these numbers at every embedding.
Conjugating the original place transports the same flag to
any other place needed for this statement.  By (CX.16) the
whole representation is de Rham; exactness of its de Rham
filtration now makes its Hodge--Tate multiset precisely
$(0,-1,-e,-e-1)$.  We have not inferred that an arbitrary
extension of de Rham characters is de Rham.
The automorphic Galois theorem identifies these numbers with
$\lambda_i+4-i$.  Sorting them gives exactly (CX.4), at both
embeddings, and hence (CX.3).  Finally its away-$\ell$
local-global compatibility transports through the actual
isomorphism (CX.16).  The decomposition groups, inertia
groups, and coefficient embedding remain those with which
we started.
\end{proof}

\begin{proposition}[CX5: the actual conductor bound and the level $M^7$]
\label{claim:CX5}
In CX4, the same cuspidal representation $\pi$ has a nonzero
$K_K(M^7)$-fixed vector.  If $p\nmid M$ and $\ell\ne p$, then
at every place $v\mid p$ the original $r_K$ is unramified and
its geometric Frobenius polynomial is $A_v(\pi)$.  This holds
at every original place in $\Omega$ above $v$ and every
corresponding Frobenius.
\end{proposition}

\begin{proof}
The second disjunct of \texttt{SafeRA} says that no two
specified weights differ by one.  It is false because
$m_1=m_2+1$ in one-based indexing.  Thus its first disjunct
holds and
\[
                              \ell\mid M.                    \tag{CX.17}
\]
This elementary observation is what permits a polynomial
level enlargement at $\ell$.  There is no exceptional
uncontrolled $\ell$-factor in the catalogue.

For a place $v$ of $K$, write $a_v=v(M)$, using the valuation
with $v(\mathfrak p_v)=1$.  Every finite character obtained by
restricting one of the original \texttt{KillsMu}$\,M$
characters to $G_{K_v}$ has conductor exponent at most $a_v$.
Here is the ramification check, including ramified $K_v$.
Local reciprocity and its norm compatibility identify the
cyclotomic action of a unit with the inverse of its norm in
the relevant cyclotomic quotient.  If $a_v>0$, a unit in
$1+M\mathcal O_{K_v}$ acts trivially, since
\[
 N_{K_v/\mathbb Q_b}(1+M x)\equiv1\pmod{M\mathbb Z_b},
                                                               \tag{CX.18}
\]
where $b$ is the residue characteristic.  To prove the
congruence, express multiplication by $1+Mx$ in an integral
$\mathbb Z_b$-basis of $\mathcal O_{K_v}$; its matrix is
$1+M A$, and its determinant is $1$ modulo $M$.  The ideal
$M\mathcal O_{K_v}$ is $\mathfrak p_v^{a_v}$.
If $a_v=0$, all $M$-th roots of unity lie in the maximal
unramified extension at $b$, by the separability of
$X^M-1$ in the residue field and Hensel lifting.  The character
is then unramified.  These arguments prove the stated
conductor bound in every case.

At $v\nmid\ell$, \texttt{LevelAway} diagonalizes the
restriction of $r_K$ to inertia into four such finite
characters.  Its inertia image is finite, so its monodromy
operator is zero, and its Artin conductor is at most $4a_v$.
Full local-global compatibility in CX4 and conductor
compatibility of local Langlands give
$a(\pi_v)\le4a_v$.  In particular, when $a_v=0$, the
Weil--Deligne parameter is unramified with zero monodromy;
$\pi_v$ is spherical.

At $v\mid\ell$, (CX.16) is de Rham and potentially
semistable.  The original local flag and exactness of the
potentially semistable functor identify the semisimplified
Weil representation with four characters.  Their inertial
parts are exactly the finite characters just bounded;
cyclotomic and unramified twists contribute no inertia to
their Weil--Deligne characters.  Semisimplified local-global
compatibility identifies this semisimple Weil representation
with that of $\pi_v\otimes|\det|^{-3/2}$.  It does not
identify the monodromy operators, and we do not need it to.
For any four-dimensional Weil--Deligne parameter $(V,\rho,N)$,
\[
 \begin{aligned}
 a(\rho,N)
  &=a(\rho)+\dim V^I-\dim(\ker N)^I,\\
  &\le a(\rho)+\operatorname{rank}N
   \le a(\rho)+3.
 \end{aligned}                                                \tag{CX.19}
\]
The first equality is the definition of its conductor; the
last inequality uses nilpotence of $N$ in dimension four.
Conductor is unchanged by an unramified twist.  Since
$a_v\ge1$ at $v\mid\ell$ by (CX.17), we obtain
\[
                    a(\pi_v)\le4a_v+3\le7a_v.                \tag{CX.20}
\]
Thus (CX.20) holds at every finite place, with zero conductor
outside $M$.

Every $\pi_v$ is generic.  The local newvector theorem gives
a nonzero vector fixed by $K_1(\mathfrak p_v^{a(\pi_v)})$.
The principal subgroup congruent to $1$ modulo
$\mathfrak p_v^{7a_v}$ is contained in that group.  At
$a_v=0$ use the spherical vector.  Their restricted tensor
product is a nonzero vector fixed by the actual adelic
principal subgroup $K_K(M^7)$.  No character conductor or
coefficient prime has been inserted into this level after
the fact.

For $p\nmid M$ and $\ell\ne p$, the preceding away-$\ell$
argument gives unramifiedness and the exact polynomial
$A_v(\pi)$.  Different lifts of geometric Frobenius differ
by inertia and hence have identical matrices under $r_K$.
Different places of $\Omega$ above $v$ are conjugate by
$G_K$, and conjugating the place conjugates the Frobenius
matrix.  This proves all of the place and Frobenius
quantifiers in the proposition.
\end{proof}

\subsection{Fixed quadratic fields and the finite ambiguity of restriction}

\begin{proposition}[CX6: eighteen fixed fields and the two-extension bound]
\label{claim:CX6}
There are eighteen imaginary quadratic fields $K_i$, fixed
once and for all, with pairwise coprime discriminants.
For every prime $p$, at least seventeen are unramified at $p$.
For every irreducible four-dimensional representation $r$
over an algebraically closed characteristic-zero field, at
least one of those seventeen quadratic restrictions is
irreducible.  Once an irreducible restriction to a fixed
$K_i$ is specified, it has at most two extensions to
$G_{\mathbb Q}$ up to isomorphism.
\end{proposition}

\begin{proof}
Choose eighteen distinct primes $d_i\equiv3\pmod4$ and put
$K_i=\mathbb Q(\sqrt{-d_i})$.  There are enough such primes
by the elementary Euclidean argument: if all had occurred
in a finite list, $4\prod_i d_i-1$ would have a new prime
factor congruent to $3$ modulo $4$.  Indeed it is $3$ modulo
$4$, some prime factor with odd multiplicity has that
congruence, and no old $d_i$ divides it.  The quadratic
integer-basis calculation gives
$\operatorname{disc}K_i=-d_i$.  Thus a prime divides at
most one discriminant.  Each remaining residue degree
$f_i$ over $p$ is either one or two.

For a representation $V$ of dimension $n$ and every distinct
self-twist character $\delta$, choose an invertible
intertwiner $T_\delta:V\simeq V\otimes\delta$.
Under conjugation on $\operatorname{End}(V)$ these are
nonzero eigenvectors for distinct one-dimensional
characters.  They are linearly independent.  One direct
proof takes a nontrivial relation of minimal length and
applies a group element distinguishing two characters;
subtracting the appropriate scalar multiple of the
original relation gives a shorter relation, a contradiction.
Consequently there are at most $n^2$ self twists, hence at
most sixteen in dimension four.  The seventeen quadratic
characters belonging to the selected fields are distinct.
One, say $\delta_i$, is not a self twist.

Write $H=G_{K_i}\triangleleft G_{\mathbb Q}$ and choose
$t\notin H$.  If a nonzero proper $H$-stable subspace $W$
existed, then $W\cap tW$ and $W+tW$ would be stable under
the whole group, since $t^2\in H$.  Irreducibility forces
$W\cap tW=0$ and $W+tW=V$.  The operator which is $+1$
on $W$ and $-1$ on $tW$ commutes with $H$ and anticommutes
with the other coset.  It is therefore an invertible
intertwiner between $r$ and $r\otimes\delta_i$, contrary
to the choice of $i$.  Thus $r|_H$ is irreducible.  This
argument requires no unproved semisimplicity assertion
about an arbitrary restriction.

Finally suppose $r|_H\simeq r'|_H$ is irreducible.  Conjugate
$r'$ so that the two restrictions agree.  The operator
$r(t)^{-1}r'(t)$ commutes with $r(H)$ and is scalar by
Schur's lemma.  An elementary proof of Schur's lemma takes
an eigenvalue of this operator: its nonzero eigenspace is
$H$-stable and must be all of $V$.  Let the scalar be $a$.
Since $t^2\in H$, it satisfies $a^2=1$.  Hence $a=1$ or
$a=-1$, and the two possible extensions are $r$ and
$r\otimes\delta_i$.  The argument works in the original
algebraic closure after any embedding of the fixed $K_i$.
\end{proof}

\subsection{The common catalogue and the literal integer endpoint}

\begin{proposition}[CX7: complete fixed-field catalogues before the coefficient prime]
\label{claim:CX7}
Fix one imaginary quadratic field $K$.  For $e\ge2$, $M\ge1$,
a prime $p\nmid M\operatorname{disc}K$ and a place $v\mid p$,
there is a finite labelled catalogue $\mathcal D_K(e,M,v)$
whose four-element multisets satisfy the literal
\texttt{EigData} conditions at $(q_v,m,e+1)$.  Its cardinality
is at most
\[
                    c_K M^{224}(e+2)^{32}.                   \tag{CX.21}
\]
For every odd $\ell\ne p$, every original $\Omega$ and
$\iota$, and every irreducible restriction to $K$ of an
original \texttt{AdmRAlt} representation, there is a label
with its exact arithmetic Frobenius multiset at $v$.
Equal labels give isomorphic restrictions.  The catalogue
is independent of all those later choices.
\end{proposition}

\begin{proof}
Take the label set to be $\mathcal S_K(e,M^7)$ from CX2 and
attach $\beta_v(\pi)$ to the label $[\pi]$.  CX1 gives
(CX.21); CX2 gives the permutation under every automorphism
of $\mathbb C$; CX3 gives the four-entry cardinality,
\texttt{Win}, and both universally quantified integral
conditions.  Thus this is the literal finite data structure,
even though its proof used a stronger pure absolute-value
bound.  The level appearing in its label set is explicitly
$M^7$.

Given one of the later actual restrictions, CX4 and CX5
produce a label in that already chosen set.  CX5 identifies
its geometric Frobenius multiset with $A_v(\pi)$, so (CX.2)
identifies its arithmetic multiset with the assigned
$\beta_v(\pi)$.  Equal labels give the same automorphic
Galois representation and hence isomorphic restrictions by
(CX.16).  Equivalently, equality of full global Satake
systems would suffice, by Chebotarev and semisimplicity.
We keep the full labels and therefore never replace this
argument by an inference from one Frobenius polynomial.
The label set was formed using only $K,e,M,v$; no choice
made to classify an individual restriction changes it or
the fixed constant $c_K$.
\end{proof}

\begin{proposition}[CX8: the entire SP9 catalogue with uniform quantifiers]
\label{claim:CX8}
There is an integer $c_*\ge1$ such that, for every $e\ge2$,
$M\ge1$ and prime $p\ge5$ with $p\nmid M$, the full
catalogue required before SP9 exists with
\[
                   A_*=224,
 \qquad |\mathcal D|\le c_*[M(e+2)]^{224}.                    \tag{CX.22}
\]
It is chosen before $\ell$, $\Omega$, $\iota$, the places
and Frobenius elements in $\Omega$, and the representation
family.  It classifies every original \texttt{AdmRAlt}
representation with weights (CX.1), has algebraic-integer
entries, and is permuted with multiplicity by every
automorphism of $\mathbb C$.  In fact every entry has
absolute value $p^{(e+1)/2}$, so its trace satisfies SP9's
bound.
\end{proposition}

\begin{proof}
Use the eighteen fields fixed in CX6.  For this $p$, let
$I_p$ consist of the indices with $p\nmid\operatorname{disc}K_i$.
Choose one place $v_i$ of $K_i$ over $p$, and let
$f_i\in\{1,2\}$ be its residue degree.  Apply CX7 to get
a catalogue $\mathcal D_i$ at $q_i=p^{f_i}$.
All these choices precede $\ell$, $\Omega$ and $\iota$.

Define $\mathcal D$ to have the following labels:
\[
 (i,s,B,\epsilon),\quad
 i\in I_p,\ s\in\mathcal D_i,\quad
 B\text{ a four-element multiset with }B^{f_i}=\beta_{i,s},
 \quad \epsilon\in\{0,1\}.                                  \tag{CX.23}
\]
The multiset attached to this label is $B$ itself.
Here $B^{f_i}$ means raising every entry to that power,
with multiplicity.  Enumerate the four entries of
$\beta_{i,s}$ for the purpose of counting.  There are at
most $f_i^4\le16$ choices of their roots and therefore at
most sixteen possible root multisets.  Hence
\[
 \begin{aligned}
 |\mathcal D|
   &\le32\sum_{i\in I_p}|\mathcal D_i|\\
   &\le32\left(\sum_{i=1}^{18}c_{K_i}\right)
                   M^{224}(e+2)^{32}
    \le c_*[M(e+2)]^{224},
 \end{aligned}                                                \tag{CX.24}
\]
where $c_*=32\sum_i c_{K_i}$ is fixed once and for all.

For an entry $b\in B$, $b^{f_i}$ is an algebraic integer,
so $b$ is integral over the algebraic integers and hence
over $\mathbb Z$.  Moreover
$(p^{e+1}b^{-1})^{f_i}=q_i^{e+1}(b^{f_i})^{-1}$ is
integral by CX3, proving the inverse integral condition
as well.  CX3 also gives
$|b|^{f_i}=q_i^{(e+1)/2}$, so
$|b|=p^{(e+1)/2}$.  An automorphism $\gamma$ permutes
the labels by
\[
        (i,s,B,\epsilon)\longmapsto
          (i,\gamma s,\gamma B,\epsilon).                    \tag{CX.25}
\]
The power equation in (CX.23) is preserved by (CX.8), and
the inverse permutation uses $\gamma^{-1}$.  This proves
the permutation with all multiplicities retained, even
when distinct labels give the same multiset.  Finally
\[
       \left|\sum B\right|
          \le4p^{(e+1)/2}
          \le(3\sqrt p+1)p^{(e+1)/2}.                         \tag{CX.26}
\]

It remains to prove the uniform classification, which is
stronger than the finiteness of (CX.23).  Fix the later
$\ell,\Omega,\iota$.  Choose a place $B_0$ of $\Omega$
above $p$.  Embed each fixed $K_i$ in $\Omega$ so that
$B_0$ restricts to the selected $v_i$.  Such a choice is
possible: at an inert prime there is only one place, and
at a split prime the two embeddings interchange the two
places.  Changing this embedding only conjugates the
restriction of a representation of $G_{\mathbb Q}$ and
does not change any of the counting bounds.

For an actual admissible $r$, CX6 supplies an index
$i\in I_p$ with irreducible restriction.  Choose the
least such index.  CX7 attaches its restriction to a
label $s\in\mathcal D_i$.  At $p$ the original $r$ is
unramified by \texttt{LevelAway}, since $p\nmid M$ and
$p\ne\ell$.  Let $B(r)$ be the complex eigenmultiset
of an arithmetic Frobenius $\sigma$ at $B_0$.  In the
unramified quadratic extension, $\sigma^{f_i}$ fixes
$K_i$ and is an arithmetic Frobenius over $v_i$.
Therefore
\[
                     B(r)^{f_i}=\beta_{i,s}.                  \tag{CX.27}
\]
This is an identity of characteristic-polynomial roots
with multiplicity, so no diagonalizability of the
Frobenius matrix is required.

For each fixed pair $(i,s)$, the admitted restrictions
are isomorphic by CX7.  By CX6 they arise from at most
two isomorphism classes of the original representation
of $G_{\mathbb Q}$.  Choose distinct bits for those
classes.  Sending $r$ to
$(i,s,B(r),\epsilon(r))$ now gives a classification map
into the fixed set (CX.23).  Equal labels imply the same
original isomorphism class.  The choice of bits is made
only for the classification map; it adds no labels and
does not change the already fixed catalogue.

Different arithmetic Frobenius lifts at $B_0$ differ by
inertia, which acts trivially.  Every other original
place above $p$ is a $G_{\mathbb Q}$-conjugate of $B_0$,
so its arithmetic Frobenius has a conjugate matrix.
Consequently the assigned label satisfies
\[
       \iota\bigl(\operatorname{tr}r(\sigma')\bigr)
                   =\sum B(r)                               \tag{CX.28}
\]
for every place and every Frobenius in the source
quantifiers, not just the auxiliary $B_0$ used in the
construction.  We have thus proved exactly all three
SP9 catalogue conditions, in their original order of
quantification.
\end{proof}

\begin{proposition}[CX9: literal \texttt{NumberSp4Alt}]
\label{claim:CX9}
The source predicate \texttt{NumberSp4Alt} holds
relative only to the classical theorems specified above.
One may take
\[
                  A=225,\qquad C=c_*2^{225}.                  \tag{CX.29}
\]
In particular, after $e\ge2$, $M\ge1$, $p\ge5$ prime
with $p\nmid M$, $k\ge1$, and three $k$-th roots of
unity $z_1,z_2,z_3$ are fixed, there is a nonzero
integer $N_*$ with
\[
 \log|N_*|\le C\varphi(k)[M(e+1)]^{225}\log p.                 \tag{CX.30}
\]
It is chosen before every odd $\ell\ne p$, every original
$\Omega,\iota$, every place and arithmetic Frobenius
$\sigma$ above $p$, and every finite pairwise
nonisomorphic family $\mathcal X$ of the original
\texttt{AdmRAlt} representations.  For all of them,
every difference below is nonzero and
\[
 |N_*|_\ell\le
 \prod_{r\in\mathcal X}
 \left|\operatorname{tr}r(\sigma)-
 \iota^{-1}\bigl(p^{e+1}z_1+p^e z_2+pz_3+1\bigr)\right|_\ell.
                                                               \tag{CX.31}
\]
\end{proposition}

\begin{proof}
CX8 supplies the exact antecedent of SP9, with $A_*=224$
and the same $c_*$ for all parameters.  Apply its proved
integral-product construction.  Explicitly its factors
are indexed by the labels in (CX.23) and by
$u\in(\mathbb Z/k)^\times$:
\[
 N_* =\prod_{s,u}
 \left(\sum\beta_s-
     p^{e+1}z_1^u-p^e z_2^u-pz_3^u-1\right).                 \tag{CX.32}
\]
CX8 gives integrality of every factor and a permutation
of the label multiset under every complex automorphism;
the same automorphism acts on the three roots of unity
by one common $u$.  The product is consequently a
rational algebraic integer.  The trace separation and
the bound on its nonzero factors are exactly the
elementary inequalities proved in SP9, with (CX.26)
as their hypothesis.  Its constants give (CX.29) and
(CX.30).  For a particular family, equal labels imply
isomorphism, so each member selects a different factor
of (CX.32).  All unselected factors have $\ell$-adic
norm at most one.  Equation (CX.28) identifies the
selected factors with exactly those in (CX.31), with
the original arithmetic Frobenius and coefficient
isomorphism.  This proves the normalized valuation
inequality and all the stated universal quantifiers.
The integer itself was formed before any of those
later choices.  This is the literal source endpoint.
\end{proof}

\subsection{What has closed, and what the closure does not assert}

CX1--CX8 remove the actual uniform catalogue premise of
SP9.  The bound is uniform because the eighteen number
fields, their cell counts and their coefficient bounds
are fixed before $e,M,p$; the finite catalogue is then
fixed before $\ell,\Omega,\iota$ or a representation
family.  CX4--CX5 use the same actual restriction of
the original Galois representation.  The passage from
$M$ to $M^7$ takes place only in the auxiliary
automorphic label set.  CX6--CX8 return to the original
representation, with its original level, through a
proved two-element ambiguity and the exact Frobenius
power identity.

The full source predicate \texttt{Hypotheses'} has not been promoted.
The general record \texttt{Literature\allowbreak RankN\allowbreak Hypothesis} also remains
separate.  Neither have unchanged-level
\texttt{LevelPNAlt} and all-rank \texttt{CountQAlt}.
In particular, this chapter
does not construct \texttt{IsAutQA} for a new Galois
representation and does not prove
\texttt{SymplecticFourImageDense}.  It proves the
literal universal counting theorem whose inputs are
already \texttt{AdmRAlt}, and closes precisely the
catalogue edge needed by the paired-count chain.

\paragraph{Primary references.}
The cohomological and rationality results are in
L.~Clozel, \emph{Motifs et formes automorphes}, in
\emph{Automorphic Forms, Shimura Varieties, and
$L$-functions}, vol.~I (1990), pp.~77--159,
Theorem~3.13 and Lemma~3.14;
\href{https://www.jmilne.org/math/Books/AA1988a.pdf}
{the editor's copy} is the primary locator.
For an explicit modern statement with the same
$C$-algebraic convention, see S.~Patrikis,
\href{https://people.math.osu.edu/patrikis.1/variationsrevision.pdf}
{\emph{Variations on a theorem of Tate}},
Theorem~2.4.5; the cohomological specialization is also
spelled out by S.~W.~Shin and N.~Templier,
\href{https://math.berkeley.edu/~swshin/FieldRat.pdf}
{\emph{On fields of rationality for automorphic
representations}}, Lemma~2.12 and Proposition~2.15.
The exact injection used here is A.~Borel,
J.-P.~Labesse and J.~Schwermer,
\href{https://www.numdam.org/item/CM_1996__102_1_1_0.pdf}
{\emph{On the cuspidal cohomology of $S$-arithmetic
subgroups}}, Theorem~5.2.
The finite compactification is A.~Borel and J.-P.~Serre,
\emph{Corners and arithmetic groups},
\emph{Comment.~Math.~Helv.} 48 (1973), 436--491,
Theorem~9.3 and the manifold-with-corners construction,
\href{https://doi.org/10.1007/BF02566134}{original article}.

For the Galois, Hodge and purity assertions see
T.~Barnet-Lamb, T.~Gee, D.~Geraghty and R.~Taylor,
\href{https://www.ma.imperial.ac.uk/~tsg/Index_files/potentialautomorphy.pdf}
{\emph{Potential automorphy and change of weight}},
Theorem~2.1.1, and
\href{https://arxiv.org/pdf/1105.2242}
{\emph{Local-global compatibility for $l=p$, II}},
Theorem~A.  The finite-coefficient and local
normalizations are checked in CX3--CX5; only
semisimplification at $\ell$ is used in (CX.19).
The filtered-module fact is P.~Colmez and J.-M.~Fontaine,
\href{https://webusers.imj-prg.fr/~pierre.colmez/CF.pdf}
{\emph{Construction des repr\'esentations $p$-adiques
semi-stables}}, Proposition~4.2 and the corollary to
Theorem~4.3.
For the implication from de Rham to potentially semistable see
L.~Berger,
\href{https://arxiv.org/pdf/math/0102179}
{\emph{Repr\'esentations $p$-adiques et \'equations
diff\'erentielles}}, Theorems~0.6 and~5.19.
For conductors see P.~Scholze,
\href{https://people.mpim-bonn.mpg.de/scholze/LocalLanglands.pdf}
{\emph{The local Langlands correspondence for
$\mathrm{GL}_n$ over $p$-adic fields}}, Theorem~14.1,
and H.~Jacquet, I.~I.~Piatetski-Shapiro and J.~Shalika,
\href{https://lesesvre.perso.math.cnrs.fr/newforms-references/jpss.pdf}
{\emph{Conducteur des repr\'esentations du groupe
lin\'eaire}}, Theorem~5.1.
Local norm compatibility and cyclotomic reciprocity are
used in their classical finite form; convenient primary
author references are J.~S.~Milne,
\href{https://www.jmilne.org/math/CourseNotes/CFT.pdf}
{\emph{Class Field Theory}}, local reciprocity and
global reciprocity chapters, and
\href{https://www.jmilne.org/math/CourseNotes/ANT.pdf}
{\emph{Algebraic Number Theory}}, Theorem~8.31
(Chebotarev).  Strong multiplicity one and genericity
are the standard general-linear cuspidal theorems;
J.~W.~Cogdell's
\href{https://people.math.osu.edu/cogdell.1/fields-www.pdf}
{\emph{Lectures on $L$-functions, converse theorems,
and functoriality for $\mathrm{GL}_n$}}, Theorem~4.3
and Corollary~4.1.2, gives the precise general-linear formulations.

